\chapter*{Conclusion}
\addcontentsline{toc}{chapter}{Conclusion}

Ce stage de perfectionnement effectué au sein de la société \textbf{SCPC Comptoir Hammami} a été une expérience professionnelle et humaine d'une grande richesse. Il m'a offert l'opportunité de m'immerger dans le monde de l'entreprise et de comprendre les enjeux liés à la digitalisation des processus métiers.

\section*{Bilan du travail effectué}
L'objectif principal de ce stage était la conception et le développement de l'application "Central de Réclamation". Je peux affirmer que cet objectif a été atteint. L'application développée permet aujourd'hui :
\begin{itemize}
    \item Aux clients de saisir leurs réclamations en ligne de manière autonome.
    \item Au service qualité de centraliser toutes les demandes sur une plateforme unique.
    \item De suivre l'état d'avancement de chaque dossier en temps réel.
    \item De sécuriser les données qui étaient auparavant dispersées sur des supports papier ou Excel.
\end{itemize}

\section*{Apports du stage}
Sur le plan personnel et technique, ce stage m'a permis de :
\begin{itemize}
    \item \textbf{Consolider mes compétences techniques :} J'ai pu approfondir ma maîtrise du langage PHP et de la conception de bases de données relationnelles MySQL. J'ai également appris à structurer mon code selon une architecture MVC pour le rendre maintenable.
    \item \textbf{Découvrir le travail en équipe :} La collaboration avec le Responsable IT et les échanges avec le service Qualité m'ont appris l'importance de la communication et de l'écoute active pour bien cerner les besoins.
    \item \textbf{Gérer un projet :} J'ai appris à respecter un planning, à prioriser les tâches et à faire face aux imprévus techniques.
\end{itemize}

\section*{Perspectives d'amélioration}
Bien que fonctionnelle, l'application pourrait être enrichie par de nouvelles fonctionnalités :
\begin{itemize}
    \item \textbf{Module de statistiques avancé :} Génération automatique de graphiques (Camemberts, Histogrammes) pour analyser les types de défauts les plus fréquents.
    \item \textbf{Application Mobile :} Développement d'une version mobile native (Android/iOS) pour permettre aux clients d'envoyer des photos directement depuis leur smartphone sur le chantier.
    \item \textbf{Chatbot :} Intégration d'un assistant virtuel pour répondre aux questions fréquentes des clients avant qu'ils ne créent un ticket.
\end{itemize}

Pour conclure, cette expérience confirme ma volonté de poursuivre mes études dans le domaine du développement web et de l'ingénierie logicielle, afin de relever des défis encore plus complexes à l'avenir.
